\documentclass[a4paper,11pt]{article}
\usepackage[utf8]{inputenc}
\usepackage[T1]{fontenc}
\usepackage[ngerman]{babel}
\usepackage[margin=2.2cm, top=2.5cm, bottom=2.5cm]{geometry}
\usepackage{enumitem}
\usepackage{booktabs}
\usepackage{longtable}
\usepackage{tabularx}
\usepackage{xcolor}
\usepackage{fancyhdr}
\usepackage{titlesec}
\usepackage{tcolorbox}
\usepackage{microtype}
\usepackage{hyperref}
\usepackage{listings}
\usepackage{graphicx}
\usepackage{parskip}
\usepackage{array}
\usepackage{mdframed}

% ── Farben ──────────────────────────────────────────────────────────
\definecolor{beezubi}{RGB}{0, 90, 160}
\definecolor{hinweis}{RGB}{255, 243, 205}
\definecolor{hinweisrand}{RGB}{200, 160, 0}
\definecolor{achtung}{RGB}{255, 220, 220}
\definecolor{achtungrand}{RGB}{180, 0, 0}
\definecolor{erfolg}{RGB}{220, 245, 220}
\definecolor{erfolgrand}{RGB}{0, 140, 0}
\definecolor{codebg}{RGB}{245, 245, 245}
\definecolor{codeframe}{RGB}{180, 180, 180}
\definecolor{grau}{RGB}{100, 100, 100}

% ── Kopf- und Fusszeile ─────────────────────────────────────────────
\pagestyle{fancy}
\fancyhf{}
\fancyhead[L]{\small\color{beezubi}\textbf{BeeZubi Lernwelt GmbH}}
\fancyhead[R]{\small\color{grau}Benutzeranleitung -- CSV-Übersetzungsprogramm}
\fancyfoot[C]{\small\thepage}
\renewcommand{\headrulewidth}{0.4pt}

% ── Überschriften ───────────────────────────────────────────────────
\titleformat{\section}{\large\bfseries\color{beezubi}}{}{0em}{}[\vspace{-0.3em}\rule{\linewidth}{0.4pt}]
\titleformat{\subsection}{\normalsize\bfseries\color{beezubi}}{}{0em}{}
\titlespacing*{\section}{0pt}{1.5ex}{0.8ex}
\titlespacing*{\subsection}{0pt}{1ex}{0.4ex}

% ── Listings (Code-Blöcke) ──────────────────────────────────────────
\lstset{
  backgroundcolor=\color{codebg},
  frame=single,
  rulecolor=\color{codeframe},
  basicstyle=\small\ttfamily,
  breaklines=true,
  breakatwhitespace=false,
  keepspaces=true,
  columns=fullflexible,
  xleftmargin=6pt,
  xrightmargin=6pt,
  aboveskip=6pt,
  belowskip=6pt,
}

% ── Boxen ───────────────────────────────────────────────────────────
\newenvironment{hinweisbox}{%
  \begin{mdframed}[backgroundcolor=hinweis,linecolor=hinweisrand,linewidth=1pt,
    innerleftmargin=8pt,innerrightmargin=8pt,innertopmargin=6pt,innerbottommargin=6pt]
  \small}{%
  \end{mdframed}}

\newenvironment{achtungbox}{%
  \begin{mdframed}[backgroundcolor=achtung,linecolor=achtungrand,linewidth=1pt,
    innerleftmargin=8pt,innerrightmargin=8pt,innertopmargin=6pt,innerbottommargin=6pt]
  \small}{%
  \end{mdframed}}

\newenvironment{erfolgbox}{%
  \begin{mdframed}[backgroundcolor=erfolg,linecolor=erfolgrand,linewidth=1pt,
    innerleftmargin=8pt,innerrightmargin=8pt,innertopmargin=6pt,innerbottommargin=6pt]
  \small}{%
  \end{mdframed}}

\setlength{\parindent}{0pt}

% ────────────────────────────────────────────────────────────────────
\begin{document}

% ── Titelseite ──────────────────────────────────────────────────────
\begin{center}
  \vspace*{1cm}
  {\LARGE\bfseries\color{beezubi} Benutzeranleitung}\\[0.4em]
  {\Large CSV-Übersetzungsprogramm}\\[0.6em]
  {\large\color{grau}\texttt{csv\_translator.py}}\\[1.2em]
  \rule{0.6\linewidth}{0.4pt}\\[0.6em]
  {\normalsize BeeZubi Lernwelt GmbH \quad|\quad Stand: März 2026}\\[0.3em]
  {\small\color{grau}Für Anwender ohne Programmierkenntnisse}
  \vspace{0.8cm}
\end{center}

% ── Überblick ────────────────────────────────────────────────────────
\section{Was macht dieses Programm?}

Das Programm übersetzt Prüfungsfragen automatisch aus einer deutschen CSV-Datei in
mehrere Zielsprachen (z.\,B. Arabisch, Türkisch, Russisch, Spanisch, Polnisch,
Englisch, Rumänisch). Es nutzt dazu die OpenAI-API (künstliche Intelligenz).

\textbf{Eingabe:}
\begin{itemize}[itemsep=2pt]
  \item Eine CSV-Datei mit deutschen Prüfungsfragen (z.\,B. \texttt{Bankkaufleute\_260202.csv})
  \item Eine Word-Datei (.docx) mit den Übersetzungsregeln (sog. Megaprompt,
        z.\,B. \texttt{Megaprompt\_26\_03\_01\_Bank\_Uebersetzen.docx})
\end{itemize}

\textbf{Ausgabe:}
\begin{itemize}[itemsep=2pt]
  \item Pro Zielsprache eine übersetzte CSV-Datei im Ausgabeordner
        \texttt{out\_translated\_backofff/}
  \item Eine Protokolldatei mit Statistiken über jeden Übersetzungslauf
\end{itemize}

% ── Schritt 1: Python installieren ───────────────────────────────────
\section{Schritt 1 — Python installieren}

Das Programm benötigt \textbf{Python 3.10 oder neuer}. Falls Python bereits
installiert ist, können Sie diesen Schritt überspringen.

\subsection{Überprüfen, ob Python bereits vorhanden ist}

Öffnen Sie \textbf{PowerShell} (Windows-Taste drücken, \texttt{PowerShell}
eingeben, Enter drücken) und geben Sie ein:

\begin{lstlisting}
python --version
\end{lstlisting}

\begin{hinweisbox}
\textbf{Ergebnis OK:} Wenn eine Versionsnummer erscheint (z.\,B.
\texttt{Python 3.12.0}), ist Python bereits installiert — weiter mit Schritt~2.\\[4pt]
\textbf{Ergebnis Fehler:} Wenn eine Fehlermeldung erscheint, muss Python
installiert werden.
\end{hinweisbox}

\subsection{Python herunterladen und installieren}

\begin{enumerate}[leftmargin=1.5em, itemsep=4pt]
  \item Rufen Sie folgende Adresse in Ihrem Browser auf:\\
        \url{https://www.python.org/downloads/}
  \item Klicken Sie auf den gelben Button \textbf{„Download Python 3.x.x"}.
  \item Starten Sie die heruntergeladene Installationsdatei (\texttt{.exe}).
  \item \textbf{Wichtig:} Aktivieren Sie ganz unten im Installer-Fenster das
        Häkchen \textbf{„Add Python to PATH"}, bevor Sie auf „Install Now" klicken.
  \item Warten Sie, bis die Installation abgeschlossen ist.
  \item Öffnen Sie eine neue PowerShell und prüfen Sie mit \texttt{python --version},
        ob die Installation erfolgreich war.
\end{enumerate}

% ── Schritt 2: Benötigte Pakete installieren ─────────────────────────
\section{Schritt 2 — Benötigte Pakete installieren}

Das Programm benötigt vier Zusatzbibliotheken, die einmalig installiert werden müssen.

Öffnen Sie \textbf{PowerShell}, navigieren Sie in den Projektordner und führen Sie
folgenden Befehl aus:

\begin{lstlisting}
pip install pandas python-docx openai ftfy
\end{lstlisting}

\begin{hinweisbox}
\textbf{In den richtigen Ordner wechseln:}\\
Falls Sie sich noch nicht im Projektordner befinden, geben Sie zunächst ein:
\begin{lstlisting}
cd "D:\Arbeit\BeeZubi_Lernwelt_GmbH\Projekt_Uebersetzer"
\end{lstlisting}
(Passen Sie den Pfad ggf. an Ihren tatsächlichen Speicherort an.)
\end{hinweisbox}

Die Pakete werden automatisch aus dem Internet heruntergeladen und installiert.
Dies dauert einige Minuten. Warten Sie, bis der Befehl abgeschlossen ist
(Sie sehen wieder den PowerShell-Prompt \texttt{>}).

\textbf{Was diese Pakete sind:}
\begin{center}
\begin{tabularx}{0.92\textwidth}{@{}l X@{}}
\toprule
\textbf{Paket} & \textbf{Verwendung} \\
\midrule
\texttt{pandas} & Lesen und Schreiben von CSV-Dateien \\[3pt]
\texttt{python-docx} & Lesen der Word-Datei mit den Übersetzungsregeln \\[3pt]
\texttt{openai} & Kommunikation mit der OpenAI-API (KI-Übersetzung) \\[3pt]
\texttt{ftfy} & Automatische Korrektur von Zeichencodierungsfehlern \\
\bottomrule
\end{tabularx}
\end{center}

% ── Schritt 3: OpenAI-API-Schlüssel einrichten ───────────────────────
\section{Schritt 3 — OpenAI-API-Schlüssel einrichten}

Das Programm kommuniziert mit dem OpenAI-Dienst (ChatGPT). Dafür wird ein
persönlicher API-Schlüssel benötigt.

\begin{achtungbox}
\textbf{Der API-Schlüssel ist vertraulich!} Geben Sie ihn niemals weiter und
speichern Sie ihn nicht in öffentlichen Dateien. Die Nutzung des Schlüssels
ist kostenpflichtig — achten Sie auf Ihre Verbrauchsgrenzen im OpenAI-Dashboard.
\end{achtungbox}

\subsection{API-Schlüssel in einer .env-Datei hinterlegen (empfohlen)}

\begin{enumerate}[leftmargin=1.5em, itemsep=4pt]
  \item Erstellen Sie im Projektordner eine neue Textdatei mit dem Namen
        \textbf{\texttt{.env}} (Punkt vor dem Namen, keine Dateiendung).
  \item Öffnen Sie die Datei mit einem Texteditor (z.\,B. Editor/Notepad) und
        schreiben Sie genau folgende Zeile hinein:
\begin{lstlisting}
OPENAI_API_KEY=sk-proj-IhrSchluesselHier
\end{lstlisting}
  \item Ersetzen Sie \texttt{sk-proj-IhrSchluesselHier} durch Ihren echten Schlüssel.
  \item Speichern und schließen Sie die Datei.
\end{enumerate}

\begin{hinweisbox}
\textbf{Hinweis:} Der API-Schlüssel beginnt immer mit \texttt{sk-}. Sie finden
ihn auf der OpenAI-Webseite unter\\
\url{https://platform.openai.com/api-keys}\\
nach Anmeldung in Ihrem OpenAI-Konto.
\end{hinweisbox}

% ── Schritt 4: Eingabedateien vorbereiten ────────────────────────────
\section{Schritt 4 — Eingabedateien vorbereiten}

Stellen Sie sicher, dass folgende Dateien im Projektordner vorhanden sind:

\begin{center}
\begin{tabularx}{0.92\textwidth}{@{}l X@{}}
\toprule
\textbf{Datei} & \textbf{Beschreibung} \\
\midrule
\texttt{MeineFragendatei.csv} & Die deutsche CSV-Datei mit den Prüfungsfragen. \\[3pt]
\texttt{Megaprompt\_...docx} & Die Word-Datei mit den Übersetzungsregeln
                               (vom Projektteam bereitgestellt). \\[3pt]
\texttt{.env} & Datei mit dem API-Schlüssel (Schritt~3). \\[3pt]
\texttt{csv\_translator.py} & Das Übersetzungsprogramm selbst. \\
\bottomrule
\end{tabularx}
\end{center}

\begin{hinweisbox}
\textbf{Anforderungen an die CSV-Datei:}
\begin{itemize}[itemsep=2pt]
  \item Trennzeichen: Semikolon \texttt{;} oder Komma \texttt{,}
        (wird automatisch erkannt)
  \item Encoding: UTF-8 oder UTF-8-BOM (\texttt{utf-8-sig}) — bei Excel-Export
        \texttt{utf-8-sig} verwenden
  \item Spaltennamen in der ersten Zeile
\end{itemize}
\end{hinweisbox}

% ── Schritt 5: Programm starten ──────────────────────────────────────
\section{Schritt 5 — Programm starten}

\subsection{Standardaufruf (alle Standardsprachen)}

Öffnen Sie \textbf{PowerShell} und navigieren Sie in den Projektordner:

\begin{lstlisting}
cd "D:\Arbeit\BeeZubi_Lernwelt_GmbH\Projekt_Uebersetzer"
\end{lstlisting}

Starten Sie das Programm mit folgendem Befehl (auf Ihren Dateinamen anpassen):

\begin{lstlisting}
python csv_translator.py `
    --pdf Bankkaufleute_260202.csv `
    --prompt Megaprompt_26_03_01_Bank_Uebersetzen.docx `
    --encoding utf-8-sig
\end{lstlisting}

\begin{hinweisbox}
\textbf{Tipp:} Das Zeichen \texttt{`} (Backtick, links neben der 1-Taste) ist in
PowerShell der Zeilenfortsetzungszeichen. Sie können den Befehl auch in einer
einzigen Zeile schreiben:
\begin{lstlisting}
python csv_translator.py --pdf Bankkaufleute_260202.csv --prompt Megaprompt_26_03_01_Bank_Uebersetzen.docx --encoding utf-8-sig
\end{lstlisting}
\end{hinweisbox}

\subsection{Nur bestimmte Sprachen übersetzen}

Wenn Sie nicht alle Sprachen, sondern z.\,B. nur Englisch und Türkisch übersetzen
möchten, fügen Sie \texttt{-{}-langage} hinzu:

\begin{lstlisting}
python csv_translator.py `
    --pdf Bankkaufleute_260202.csv `
    --prompt Megaprompt_26_03_01_Bank_Uebersetzen.docx `
    --encoding utf-8-sig `
    --langage EN TR
\end{lstlisting}

\subsection{Mit angepasster Spaltenreihenfolge in der Ausgabe}

Falls die Ausgabedatei nur bestimmte Spalten in einer festgelegten Reihenfolge
enthalten soll, verwenden Sie \texttt{-{}-column-order}:

\begin{lstlisting}
python csv_translator.py `
    --pdf Bankkaufleute_260202.csv `
    --prompt Megaprompt_26_03_01_Bank_Uebersetzen.docx `
    --encoding utf-8-sig `
    --column-order "lfdNr;FrageNr;Frage;A;B;C;D;E;Richtig1;Sprache"
\end{lstlisting}

\begin{achtungbox}
\textbf{Wichtig:} Spalten, die \emph{nicht} in \texttt{-{}-column-order} aufgeführt
sind, erscheinen nicht in der Ausgabe und werden auch nicht übersetzt. Stellen Sie
sicher, dass alle gewünschten Spalten angegeben sind.
\end{achtungbox}

\subsection{Vollständiges Beispiel mit allen wichtigen Optionen}

\begin{lstlisting}
python csv_translator.py `
    --pdf Bankkaufleute_260202.csv `
    --prompt Megaprompt_26_03_01_Bank_Uebersetzen.docx `
    --encoding utf-8-sig `
    --langage AR TR RU ES PL EN RO `
    --batch-size 200 `
    --column-order "lfdNr;FrageNr;BerufNr;Beruf;LFNr;LF;AbschnNr;Abschnitt;Nr;Frage;A;B;C;D;E;Richtig1;Richtig_Text1;Schwierigkeit;Sprache"
\end{lstlisting}

% ── Schritt 6: Programmablauf verstehen ──────────────────────────────
\section{Schritt 6 — Was passiert während der Ausführung?}

Nach dem Start erscheinen Meldungen in der PowerShell. Hier ist, was diese bedeuten:

\begin{center}
\begin{tabularx}{0.96\textwidth}{@{}l X@{}}
\toprule
\textbf{Meldung} & \textbf{Bedeutung} \\
\midrule
\texttt{7149 Zeilen | 36 Batches/Sprache | 7 Sprachen} &
  Zusammenfassung: Anzahl Zeilen, Pakete pro Sprache, Anzahl Sprachen \\[4pt]
\texttt{[AR] 5.6\% | Ligne 200/7149 | Batch 0..199 OK | ETA: \textasciitilde32m} &
  Fortschrittsanzeige: Sprache, Prozent, Zeilen, ETA (geschätzte Restzeit) \\[4pt]
\texttt{[AR] JSON-Fehler — Batch wird halbiert} &
  Ein Paket konnte nicht verarbeitet werden; das Programm teilt es automatisch
  in kleinere Teile auf \\[4pt]
\texttt{[AR] RETRY 1/4: Neuübersetzung von 105 Zellen} &
  Nicht übersetzte Zellen werden automatisch nochmals übersetzt \\[4pt]
\texttt{[AR] 100.0\% | FERTIG in 36m38s} &
  Die Übersetzung für Arabisch ist abgeschlossen \\[4pt]
\texttt{[AR] Protokoll geschrieben:} &
  Die Statistiken wurden in die Protokolldatei gespeichert \\[4pt]
\texttt{Abgeschlossen.} &
  Alle Sprachen wurden verarbeitet — das Programm ist fertig \\
\bottomrule
\end{tabularx}
\end{center}

\begin{hinweisbox}
\textbf{Laufzeit:} Die Übersetzung kann je nach Dateigröße und Anzahl der Sprachen
mehrere Stunden dauern. Das Programm kann jederzeit unterbrochen und später
fortgesetzt werden — der Fortschritt wird automatisch gespeichert.
\end{hinweisbox}

% ── Schritt 7: Ergebnisse finden ─────────────────────────────────────
\section{Schritt 7 — Wo finde ich die Ergebnisse?}

Alle Ausgabedateien befinden sich im Unterordner \texttt{out\_translated\_backofff}
innerhalb des Projektordners.

\subsection{Übersetzte CSV-Dateien}

Für jede Zielsprache wird eine eigene Datei erstellt:

\begin{center}
\begin{tabularx}{0.85\textwidth}{@{}l X@{}}
\toprule
\textbf{Dateiname} & \textbf{Inhalt} \\
\midrule
\texttt{Bankkaufleute\_260202\_AR.csv} & Arabische Übersetzung \\
\texttt{Bankkaufleute\_260202\_TR.csv} & Türkische Übersetzung \\
\texttt{Bankkaufleute\_260202\_RU.csv} & Russische Übersetzung \\
\texttt{Bankkaufleute\_260202\_ES.csv} & Spanische Übersetzung \\
\texttt{Bankkaufleute\_260202\_PL.csv} & Polnische Übersetzung \\
\texttt{Bankkaufleute\_260202\_EN.csv} & Englische Übersetzung \\
\texttt{Bankkaufleute\_260202\_RO.csv} & Rumänische Übersetzung \\
\bottomrule
\end{tabularx}
\end{center}

\begin{hinweisbox}
\textbf{CSV-Dateien in Excel öffnen:} Doppelklick auf die Datei genügt in der Regel
nicht (Excel interpretiert dann kein Semikolon als Trennzeichen). Gehen Sie stattdessen
in Excel auf \textbf{Daten → Aus Text/CSV} und wählen Sie dort als Trennzeichen
\textbf{Semikolon} und als Zeichencodierung \textbf{UTF-8} aus.
\end{hinweisbox}

\subsection{Protokolldatei}

Die Datei \texttt{protokoll\_Bankkaufleute\_260202.csv} enthält für jede Sprache eine
Zeile mit folgenden Informationen:

\begin{center}
\begin{tabularx}{0.96\textwidth}{@{}l X@{}}
\toprule
\textbf{Spalte} & \textbf{Bedeutung} \\
\midrule
\multicolumn{2}{@{}l}{\textit{Abschnitt A — Technische Laufdaten}} \\[2pt]
\texttt{Datum} & Datum und Uhrzeit des Laufes \\
\texttt{Sprache} & Zielsprache (z.\,B. AR, ES) \\
\texttt{Modell} & Verwendetes KI-Modell \\
\texttt{Anzahl\_API\_Calls} & Wie oft die KI kontaktiert wurde \\
\texttt{Gesamt\_Input\_Tokens} & Gesendete Textmenge (für Kostenrechnung relevant) \\
\texttt{Gesamt\_Output\_Tokens} & Empfangene Textmenge \\
\texttt{Laufzeit} & Dauer des Laufes (z.\,B. \texttt{36m38s}) \\
\texttt{Anzahl\_Retries} & Wie viele automatische Wiederholungsrunden stattfanden \\[6pt]

\multicolumn{2}{@{}l}{\textit{Abschnitt B — Prüfergebnisse}} \\[2pt]
\texttt{Gesamtzeilen} & Anzahl Zeilen in der Eingabedatei \\
\texttt{Deutsch\_Reste\_automatisch} & Zellen, die nach dem ersten Durchlauf
                                       noch auf Deutsch waren \\
\texttt{Verdachtsfaelle\_gesamt} & Alle verdächtigen Zellen über alle Retry-Runden \\
\texttt{Manuell\_korrigiert} & \textbf{Manuell auszufüllen} nach manueller Prüfung \\
\texttt{Neu\_uebersetzt} & Wie viele Zellen durch Retry neu übersetzt wurden \\
\texttt{Endgueltige\_Deutsch\_Reste} & Verbleibende deutsche Reste nach allen
                                       Wiederholungen \\[6pt]

\multicolumn{2}{@{}l}{\textit{Abschnitt C — Automatische Empfehlungen}} \\[2pt]
\texttt{Technische\_Auffaelligkeiten} & Automatisch erkannte Auffälligkeiten \\
\texttt{Anpassungen\_am\_Prompt} & Empfehlungen zur Verbesserung der Regeldatei \\
\texttt{Anpassungen\_an\_Pipeline} & Empfehlungen zur technischen Pipeline \\
\texttt{Empfehlung\_naechster\_Lauf} & Was beim nächsten Lauf beachtet werden sollte \\
\bottomrule
\end{tabularx}
\end{center}

\subsection{Fortschrittsdateien (technisch)}

Dateien mit dem Namen \texttt{.progress\_...\_.json} im Ausgabeordner speichern
den Fortschritt jeder Sprache. Sie müssen diese Dateien nicht öffnen oder ändern.
Wenn eine Sprache fertig ist, steht \texttt{"done": true} darin.

% ── Programm unterbrechen und fortsetzen ─────────────────────────────
\section{Programm unterbrechen und fortsetzen}

Das Programm kann jederzeit mit \textbf{Strg+C} in PowerShell unterbrochen werden.

Beim nächsten Start erkennt das Programm automatisch, welche Zeilen bereits
übersetzt wurden, und setzt genau dort fort — es wird keine Arbeit doppelt gemacht.

\begin{erfolgbox}
\textbf{Beispiel:} Wenn der Computer abgestürzt ist, nachdem 3.000 von 7.149 Zeilen
übersetzt wurden, beginnt das Programm beim nächsten Start bei Zeile 3.001.
\end{erfolgbox}

% ── Fehlerbehandlung ─────────────────────────────────────────────────
\section{Häufige Probleme und Lösungen}

\begin{center}
\begin{tabularx}{0.96\textwidth}{@{}p{4cm} X@{}}
\toprule
\textbf{Problem} & \textbf{Lösung} \\
\midrule
\texttt{python} nicht gefunden &
  Python wurde nicht im PATH installiert. Python deinstallieren und neu
  installieren — dabei „Add to PATH" aktivieren. \\[6pt]

\texttt{No module named openai} &
  Pakete wurden noch nicht installiert. Befehl ausführen:\\
  \texttt{pip install pandas python-docx openai ftfy} \\[6pt]

\texttt{Bitte OPENAI\_API\_KEY setzen} &
  Die \texttt{.env}-Datei fehlt oder ist falsch geschrieben.
  Prüfen Sie, ob die Datei wirklich \texttt{.env} heißt (mit Punkt,
  ohne weitere Endung wie \texttt{.txt}). \\[6pt]

\texttt{AuthenticationError} &
  Der API-Schlüssel ist ungültig oder abgelaufen. Prüfen Sie ihn unter
  \url{https://platform.openai.com/api-keys}. \\[6pt]

Ausgabe-CSV hat komische Zeichen &
  Öffnen Sie die Datei in Excel mit Encoding UTF-8 (Daten → Aus Text/CSV). \\[6pt]

Programm läuft seit Stunden und macht nichts &
  Bei sehr großen Dateien ist eine lange Laufzeit normal. Schauen Sie auf die
  Prozentanzeige in der Konsole. Erscheint nichts mehr, können Sie das Programm
  mit Strg+C stoppen und neu starten — der Fortschritt bleibt erhalten. \\[6pt]

\texttt{PermissionError: out\_translated...} &
  Die Ausgabe-CSV ist in Excel geöffnet. Schließen Sie Excel und das Programm
  versucht automatisch erneut zu speichern. \\
\bottomrule
\end{tabularx}
\end{center}

% ── Vollständige Parameterübersicht ──────────────────────────────────
\section{Vollständige Parameterübersicht}

Alle verfügbaren Optionen des Programms im Überblick:

\begin{center}
\begin{tabularx}{0.96\textwidth}{@{}l l X@{}}
\toprule
\textbf{Parameter} & \textbf{Standard} & \textbf{Beschreibung} \\
\midrule
\texttt{-{}-pdf} & \textit{Pflicht} &
  Pfad zur Eingabe-CSV-Datei \\[4pt]
\texttt{-{}-prompt} & \textit{Pflicht} &
  Pfad zur Word-Datei mit Übersetzungsregeln \\[4pt]
\texttt{-{}-encoding} & \texttt{utf-8} &
  Zeichencodierung der CSV (\texttt{utf-8} oder \texttt{utf-8-sig} für Excel) \\[4pt]
\texttt{-{}-langage} & AR TR RU ES PL EN RO &
  Zielsprachen (durch Leerzeichen getrennt) \\[4pt]
\texttt{-{}-model} & \texttt{gpt-5-mini} &
  OpenAI-Modell (Standardmodell ist schnell und günstig) \\[4pt]
\texttt{-{}-batch-size} & \texttt{200} &
  Anzahl Zeilen pro API-Paket (10–200). Kleinere Werte = stabiler, größere = schneller \\[4pt]
\texttt{-{}-outdir} & \texttt{out\_translated\_backofff} &
  Name des Ausgabeordners \\[4pt]
\texttt{-{}-column-order} & \textit{keine} &
  Spaltenreihenfolge der Ausgabe, durch Semikolon getrennt. Nur gelistete Spalten
  werden ausgegeben \\[4pt]
\texttt{-{}-never-translate} & \textit{leer} &
  Zusätzliche Spalten, die nicht übersetzt werden sollen (kommagetrennt) \\[4pt]
\texttt{-{}-pruefung} & \textit{inaktiv} &
  Benennt Spalten automatisch um: \texttt{Zwischenprüfung}
  $\to$ \texttt{Abschlussprüfung Teil 1} \\[4pt]
\texttt{-{}-temperature} & \textit{auto} &
  Kreativitätsgrad der KI (0 = präzise, 1 = kreativ). Ohne Angabe: Modell-Standard \\[4pt]
\texttt{-{}-sep} & \textit{auto} &
  CSV-Trennzeichen, falls die automatische Erkennung nicht funktioniert \\
\bottomrule
\end{tabularx}
\end{center}

% ── Sprachcodes ──────────────────────────────────────────────────────
\section{Unterstützte Sprachen}

\begin{center}
\begin{tabular}{@{}ll@{\qquad}ll@{\qquad}ll@{\qquad}ll@{}}
\toprule
\textbf{Code} & \textbf{Sprache} &
\textbf{Code} & \textbf{Sprache} &
\textbf{Code} & \textbf{Sprache} &
\textbf{Code} & \textbf{Sprache} \\
\midrule
AR & Arabisch  & TR & Türkisch & RU & Russisch & ES & Spanisch \\
PL & Polnisch  & EN & Englisch & RO & Rumänisch & UK & Ukrainisch \\
\bottomrule
\end{tabular}
\end{center}

Weitere Sprachen können über den Parameter \texttt{-{}-langage} mit ihrem
ISO-639-1-Code angegeben werden (z.\,B. \texttt{FR} für Französisch,
\texttt{IT} für Italienisch).

% ── Checkliste ───────────────────────────────────────────────────────
\section{Checkliste vor dem ersten Start}

\begin{itemize}[label=\(\square\), itemsep=4pt, leftmargin=2em]
  \item Python 3.10 oder neuer ist installiert (\texttt{python -{}-version} prüfen)
  \item Pakete sind installiert (\texttt{pip install pandas python-docx openai ftfy})
  \item Die Datei \texttt{.env} mit dem API-Schlüssel ist im Projektordner vorhanden
  \item Die CSV-Eingabedatei ist im Projektordner vorhanden
  \item Die Megaprompt-Word-Datei ist im Projektordner vorhanden
  \item PowerShell ist auf den Projektordner gesetzt (\texttt{cd "-{}-Pfad-{}-})
  \item Der Startbefehl enthält den richtigen Dateinamen für \texttt{-{}-pdf}
        und \texttt{-{}-prompt}
\end{itemize}

% ── Kontakt ──────────────────────────────────────────────────────────
\vspace{1em}
\begin{center}
\rule{0.5\linewidth}{0.4pt}\\[0.4em]
{\small\color{grau}BeeZubi Lernwelt GmbH — Interne Dokumentation — März 2026}
\end{center}

\end{document}
